% A relaxed auxiliary norm whose integrality follows from the size bounds.
\section{Recovering the integral auxiliary Pell sequence}
\label{sec:relaxed-auxiliary}

The auxiliary norm can be weakened while preserving the index information
used by the doubled signed block. Starting with Theorem~\ref{thm:best99},
replace only E16 by
\begin{equation}\label{eq:relaxed-auxiliary-norm}
 (ic^2)^2=D(f^2-1),\qquad D=A^2-1,\qquad A=a+4.
\end{equation}
The previous norm was $f^2-1=D(ic^2)^2$. The two individual equations
are not equivalent: the new equation can admit additional rational
Pell solutions. The existing size bounds exclude the unwanted
solutions and recover the divisibility needed by the signed-index
argument. At the arithmetic level, the square $(ic^2)^2$ is now
exactly the coefficient $K=D(f^2-1)$ used by that argument, saving
one multiplication.

\subsection{A preliminary bound on the main index}

Keep the definitions
\[
 U=wn^2,\quad Y_P=sn^2,\quad a=Y_P(U+1),\quad
 A=a+4,\quad P=2UY_P^2+1,\quad J=2r+1.
\]
Before any auxiliary norm is used, the coding bootstrap supplies
\[
 n\ge64,\qquad n\le r<2n^3,\qquad U,Y_P\ge n^2,
 \qquad UY_P>r+1.
\]
The unchanged main norm and first triangular norm give positive
indices $p,t$ with
\[
 d=\chi_A(p),\quad c=\psi_A(p),\qquad
 2\tau+1=\chi_P(t),\quad k=\psi_P(t).
\]
Since $P\equiv1\pmod{UY_P}$ and $k=r+1+hUY_P$, the congruence
for $\psi_P$ gives $t\equiv r+1\pmod{UY_P}$. Thus $t\ge r+1$.
Also
\[
 P-A=UY_P(2Y_P-1)-Y_P-3>0.
\]
If $p\le r+1$, Pell growth would give
\[
 c\le(2A)^{p-1}\le(2A)^r
  <(2P-1)^r\le k,
\]
contrary to $c=kY_P+\eta>k$. Consequently
\begin{equation}\label{eq:relaxed-main-index-bound}
 p\ge r+2\ge66,\qquad c>A^6>A D^2.
\end{equation}
For the second assertion, use
$\psi_A(p)\ge(2A-1)^{p-1}$ and $D<A^2$.
This argument uses neither the auxiliary norm, the signed index
equation, an exponential identity, nor internal coefficient decoding.

\subsection{The rank argument}

\begin{lemma}\label{lem:relaxed-pell-integrality}
Let $A>1$, $D=A^2-1$, and let
$c=\psi_A(p)$, $d=\chi_A(p)$ for a positive index $p$.
Suppose $c>A D^2$ and positive integers $f,i$ satisfy
\eqref{eq:relaxed-auxiliary-norm}. Then there is a positive index $m$
such that
\[
 f=\chi_A(m),\qquad p\mid m,\qquad c\mid m,
 \qquad ic^2=D\psi_A(m).
\]
\end{lemma}
\begin{proof}
Write $D=d_0s_0^2$ with $d_0>1$ squarefree. We first justify the
use of a fundamental unit within the integer-coefficient ring
$\mathbb Z[\sqrt{d_0}]$. The known unit $A+s_0\sqrt{d_0}$ ensures
that positive solutions to $x^2-d_0y^2=1$ exist. Choose the least
positive unit
\[
 \varepsilon=F+y_0\sqrt{d_0}>1,
 \qquad F,y_0\in\mathbb Z_{>0},\qquad
 F^2-d_0y_0^2=1.
\]
Such a choice exists by minimizing its first coordinate. Every other
positive integer-coefficient unit $\beta>1$ is a power of
$\varepsilon$. Indeed, choose $j\ge0$ with
$\varepsilon^j\le\beta<\varepsilon^{j+1}$. The quotient
$\beta\varepsilon^{-j}$ has integer coefficients and norm one,
because $\varepsilon^{-1}=F-y_0\sqrt{d_0}$. If it belonged to
$(1,\varepsilon)$, both coordinates would be positive integers,
contradicting minimality. It is therefore one. No assertion about
the full ring of algebraic integers is needed.

There is accordingly an $e\ge1$ such that
\[
 A=\chi_F(e),\qquad s_0=y_0S,\qquad
 S=\psi_F(e),\qquad D=(F^2-1)S^2.
\]
Composition of the Pell solutions gives
\begin{equation}\label{eq:relaxed-composed-coordinate}
 \psi_F(ep)=Sc.
\end{equation}
Put $R_0=ic^2$. The relaxed equation gives
$R_0^2=d_0s_0^2(f^2-1)$. Thus $s_0\mid R_0$, and squarefreeness
implies $d_0\mid R_0/s_0$. The positive integer
$y=R_0/(d_0s_0)$ satisfies $f^2-d_0y^2=1$. Hence for some
$n_0>0$,
\begin{equation}\label{eq:relaxed-field-coordinate}
 f=\chi_F(n_0),\qquad
 R_0=(F^2-1)S\psi_F(n_0).
\end{equation}

We use the strong divisibility identity
\begin{equation}\label{eq:relaxed-strong-divisibility}
 \gcd(\psi_F(u),\psi_F(v))=\psi_F(\gcd(u,v)).
\end{equation}
For completeness, the Pell norm gives
$\gcd(\chi_F(j),\psi_F(j))=1$. The addition formula therefore
reduces $\gcd(\psi_F(u),\psi_F(u+v))$ to
$\gcd(\psi_F(u),\psi_F(v))$, and the Euclidean algorithm proves
\eqref{eq:relaxed-strong-divisibility}.

Let $T=(F^2-1)S$, $g=\gcd(n_0,ep)$, and
\[
 h=\gcd(\psi_F(n_0),\psi_F(ep))=\psi_F(g).
\]
Because $c^2\mid R_0=T\psi_F(n_0)$,
\[
 \begin{aligned}
 \gcd(c^2,T\psi_F(ep))
   &=\gcd(c^2,Dc)\\
   &=c\gcd(c,D)\ge c
 \end{aligned}
\]
divides $T h$. Since $T=D/S\le D$,
\begin{equation}\label{eq:relaxed-gcd-lower-bound}
 h\ge c/T\ge c/D.
\end{equation}
If $g<ep$, then $2g\le ep$, since $g$ is a proper divisor. The
duplication identity gives
\[
 2h^2<2\chi_F(g)\psi_F(g)=\psi_F(2g)
       \le\psi_F(ep)=Sc.
\]
Here $S<A$, by $D=(F^2-1)S^2$ and $F\ge2$. Thus $h^2<Ac/2$,
whereas \eqref{eq:relaxed-gcd-lower-bound} and $c>AD^2$ give
$h^2\ge c^2/D^2>Ac$, a contradiction. Therefore $ep\mid n_0$.
Writing $n_0=em$ proves $p\mid m$ and, by composition,
\begin{equation}\label{eq:relaxed-integrality-restored}
 f=\chi_A(m),\qquad R_0=D\psi_A(m).
\end{equation}

To obtain $c\mid m$, write $m=pk$. Expansion of
$(d+c\sqrt D)^k$, followed by division of the square-root
coefficient by $c$, gives
\[
 \frac{\psi_A(pk)}c\equiv k d^{k-1}\pmod c.
\]
As $c^2\mid D\psi_A(m)$ and $\gcd(d,c)=1$, this implies
$c\mid Dk$. The main Pell expansion modulo $D$ also gives
\[
 c=\psi_A(p)\equiv pA^{p-1}\pmod D,
 \qquad \gcd(c,D)=\gcd(p,D),
\]
using $\gcd(A,D)=1$. Put $g_0=\gcd(c,D)$. Then $c/g_0\mid k$
and $g_0\mid p$, so $c\mid pk=m$.
\end{proof}

The lemma recovers precisely the integral sequence and auxiliary
index divisibility needed below. It does not assert
$c^2\mid\psi_A(m)$, which is stronger than necessary and is not
an automatic consequence of the relaxed norm.

\subsection{The signed index and all positive witnesses}

By \eqref{eq:relaxed-main-index-bound},
Lemma~\ref{lem:relaxed-pell-integrality} applies before the signed
auxiliary norm. Its source polynomial remains
\[
 v=of-d^2,\qquad K=D(f^2-1)=(ic^2)^2,\qquad
 v(v+1)=K(K-1)(J+jc)^2.
\]
For $G=2K-1$, reduction modulo $f$ still gives
\[
 G\equiv-\chi_A(2),\qquad
 2v+1\equiv-\chi_A(2p)\pmod f.
\]
The conclusions $f=\chi_A(m)$ and $c\mid m$ give the comparison
bound $0<2p\le c\le m$. Also $K=(ic^2)^2$ gives
$G\equiv-1\pmod c$ directly. The signed chi step-down and parity
argument in the proof of Lemma~\ref{lem:doubled-main-index} therefore
apply unchanged: they yield $J\equiv\pm p\pmod c$, exclude
$J+p=c$, and conclude $p=J$. The old stronger divisibility of the
second Pell coordinate is not used at this step.

The exact first index, ratio estimates, common exponential witness,
fixed-index recovery through $T_L=\psi_4(L)$, and coefficient
decoding now proceed in the order proved for
Theorem~\ref{thm:best99}. All their hypotheses have been recovered
before use. This proves sufficiency of the modified system.

For necessity, start with that theorem's positive witness construction.
Its auxiliary index is the even number $m=2cJ$, and
\[
 f=\chi_A(m),\qquad i_{\mathrm{old}}=\psi_A(m)/c^2>0.
\]
Keep $f$ and set $i_{\mathrm{new}}=D i_{\mathrm{old}}>0$. Then
\[
 (i_{\mathrm{new}}c^2)^2
  =D^2\psi_A(m)^2=D(f^2-1).
\]
The source value $K=D(f^2-1)$ has not changed. Thus the same odd
$f$, the same $G$, and the positive $o,j$ constructed in
Lemma~\ref{lem:doubled-main-index} still satisfy the signed norm.
All other witnesses and fixed index components are preserved.

\subsection{One multiplication is removed}

Previously E16 and the shared coefficient $K$ required six
instructions:
\[
 \begin{gathered}
 z=ic^2,\quad z_2=z^2,\quad AE=Dz_2,\\
 f_2=f^2,\quad R_{16}=AE+1,\quad K=D\,AE.
 \end{gathered}
\]
The relaxed equation needs five:
\[
 z=ic^2,\quad K=z^2,\quad f_2=f^2,\quad
 f_-=f_2-1,\quad R_{16}=D f_-.
\]
Test $K=R_{16}$ and use the already computed square $K$ in the
signed norm. All other instructions remain unchanged.

The exact triangular correction also becomes simpler. If
\[
 F_{16}=(ic^2)^2-D(f^2-1),\quad
 K_s=D(f^2-1),\quad K_c=(ic^2)^2,
\]
then
\[
 F_{17}^{\mathrm{calc}}
 =F_{17}^{\mathrm{source}}
  -F_{16}(K_s+K_c-1)(J+jc)^2.
\]
The complete verifier checks this correction, the retained geometric
and second-exponent corrections, all source residuals and primitive
instructions, and the input domains. The Pell integrality and index
claims are proved above, independently of the symbolic checks.

\begin{theorem}\label{thm:best98-pell}
Replacing E16 by \eqref{eq:relaxed-auxiliary-norm} in the system of
Theorem~\ref{thm:best99} preserves its existential representation over
34 positive integer unknowns and 22 equality tests. Its straight-line
certificate has
\[
 \boxed{98=54\text{ multiplications}+44\text{ additions}}.
\]
Fixed numerals and supplied index parameters are free. The saving
relative to Theorem~\ref{thm:best99} is one multiplication.
\end{theorem}

The full schedule, source polynomials and exact residual identities are
recorded in the standalone checker
\begin{center}\small
\texttt{round28\_1980\_relaxed\_auxiliary\_certificate.py}
\end{center}
and its JSON receipt.
